\documentclass[a4paper,11pt]{article}
\usepackage{fontspec}
\usepackage[french]{babel}
\usepackage{amsmath,amssymb,amsthm}
\usepackage{geometry}
\usepackage{enumitem}
\usepackage{hyperref}
\usepackage{fancyhdr}
\usepackage{xcolor}
\usepackage{listings}
\usepackage{array}
\usepackage{booktabs}
\usepackage{multirow}
\usepackage{graphicx}
\usepackage{tikz}
\usetikzlibrary{shapes.geometric, arrows.meta, positioning, calc}

\geometry{margin=2cm}
\pagestyle{fancy}
\fancyhf{}
\fancyhead[C]{\textbf{STA211 — Arbres de régression et classification — Fiche récapitulative}}
\fancyfoot[C]{\thepage}
\renewcommand{\headrulewidth}{0.4pt}

\definecolor{codegreen}{rgb}{0.2,0.6,0.2}
\definecolor{codegray}{rgb}{0.5,0.5,0.5}
\definecolor{codepurple}{rgb}{0.58,0,0.82}
\definecolor{backcolour}{rgb}{0.95,0.95,0.92}

\lstdefinestyle{mystyle}{
    backgroundcolor=\color{backcolour},
    commentstyle=\color{codegreen},
    keywordstyle=\color{magenta},
    numberstyle=\tiny\color{codegray},
    stringstyle=\color{codepurple},
    basicstyle=\footnotesize\ttfamily,
    breakatwhitespace=false,
    breaklines=true,
    captionpos=b,
    keepspaces=true,
    numbers=left,
    numbersep=5pt,
    showspaces=false,
    showstringspaces=false,
    showtabs=false,
    tabsize=2
}
\lstset{style=mystyle}

\setlength{\parindent}{0pt}

\begin{document}

\begin{center}
{\LARGE\textbf{Arbres de régression et de classification}} \\
\vspace{0.2cm}
{\large \textit{CART — Classification And Regression Trees}} \\
\vspace{0.3cm}
STA211 — Apprentissage supervisé — Fiche récapitulative \\
6 juin 2026
\end{center}

\vspace{0.3cm}
\rule{\textwidth}{0.4pt}
\vspace{0.3cm}

\tableofcontents
\newpage

% ============================================================
\section{Présentation}
% ============================================================

Les \textbf{arbres CART} (Breiman et al., 1984) sont des méthodes d'apprentissage \textbf{supervisé non-paramétriques}. Ils partitionnent l'espace des variables explicatives $\mathcal{X}$ en régions et associent une prédiction simple (moyenne ou mode) à chaque région.

\paragraph{Types d'arbres :}
\begin{itemize}
    \item \textbf{Arbre de régression} : $Y$ quantitative → prédiction = moyenne des $Y$ dans la région.
    \item \textbf{Arbre de classification} : $Y$ qualitative → prédiction = modalité majoritaire dans la région.
\end{itemize}

\paragraph{Avantages :}
\begin{itemize}
    \item Interprétabilité (règles de décision lisibles)
    \item Gère naturellement les variables mixtes (quantitatives + qualitatives)
    \item Pas d'hypothèse sur la distribution des données
\end{itemize}

\paragraph{Inconvénients :}
\begin{itemize}
    \item Forte variance (instabilité) : un petit changement dans les données peut complètement changer l'arbre
    \item Biais de prédiction (souvent moins performant que forêts aléatoires / boosting)
\end{itemize}

% ============================================================
\section{Architecture d'un arbre}
% ============================================================

Un arbre est composé de :
\begin{itemize}
    \item \textbf{Racine} : premier nœud contenant tous les individus.
    \item \textbf{Nœuds internes (fils)} : division selon une variable $X_j$ et un seuil $s$.
    \item \textbf{Feuilles (nœuds terminaux)} : régions $R_1, \ldots, R_M$ de l'espace des entrées.
\end{itemize}

La valeur prédite dans une région $R_m$ est :
\[
\text{Régression : } \hat{y}_{R_m} = \frac{1}{n_m}\sum_{i \in R_m} y_i,\qquad
\text{Classification : } \hat{y}_{R_m} = \arg\max_k \hat{p}_{mk}
\]
où $\hat{p}_{mk}$ est la proportion d'individus de classe $k$ dans la région $R_m$.

Le modèle s'écrit :
\[
f(\mathbf{x}) = \sum_{m=1}^{M} \hat{y}_{R_m} \cdot \mathbf{1}_{\{\mathbf{x} \in R_m\}}
\]

\begin{center}
\begin{tikzpicture}[
    scale=0.85, transform shape,
    node distance=1.2cm and 2cm,
    dec/.style={rectangle, draw, rounded corners, minimum width=2.2cm, minimum height=0.6cm, align=center, font=\bfseries, fill=blue!10},
    leaf/.style={rectangle, draw, minimum width=2.4cm, minimum height=0.55cm, align=center, font=\small, fill=green!10},
    edge/.style={draw, -{Stealth[scale=1.2]}, thick},
    sn/.style={font=\small, sloped, above}
]

\node[dec] (r) {$X_j \le s$};

\node[dec, below left=of r] (a) {$X_k \le t$};
\node[dec, below right=of r] (b) {$X_\ell \le u$};

\node[leaf, below left=of a] (l1) {$R_1{:}\; \hat{y}_{R_1}$};
\node[leaf, below right=of a] (l2) {$R_2{:}\; \hat{y}_{R_2}$};
\node[leaf, below left=of b] (l3) {$R_3{:}\; \hat{y}_{R_3}$};
\node[leaf, below right=of b] (l4) {$R_4{:}\; \hat{y}_{R_4}$};

\draw[edge] (r) -- node[sn]{oui} (a);
\draw[edge] (r) -- node[sn]{non} (b);
\draw[edge] (a) -- node[sn]{oui} (l1);
\draw[edge] (a) -- node[sn]{non} (l2);
\draw[edge] (b) -- node[sn]{oui} (l3);
\draw[edge] (b) -- node[sn]{non} (l4);

\end{tikzpicture}
\end{center}

% ============================================================
\section{Algorithme de construction}
% ============================================================

\subsection{Principe du partitionnement récursif}

On procède \textbf{de haut en bas}, par divisions binaires successives (\textit{recursive binary splitting}) :
\begin{enumerate}
    \item Partir de la racine (tous les individus).
    \item Pour chaque nœud, chercher la variable $X_j$ et le seuil $s$ qui minimisent le \textbf{critère d'hétérogénéité}.
    \item Diviser le nœud en deux fils selon $X_j < s$ et $X_j \ge s$.
    \item Répéter sur chaque nœud fils jusqu'à un critère d'arrêt (homogénéité, taille minimale, profondeur max).
\end{enumerate}

\subsection{Critères de division}

\subsubsection{Régression (critère des moindres carrés)}
On minimise la somme des carrés résiduels (RSS) dans les deux nœuds fils :
\[
\min_{j,s} \Bigg[ \sum_{i: x_{ij} < s} (y_i - \bar{y}_{g})^2 + \sum_{i: x_{ij} \ge s} (y_i - \bar{y}_{d})^2 \Bigg]
\]
où $\bar{y}_{g}$ et $\bar{y}_{d}$ sont les moyennes dans les nœuds gauche et droit.

\subsubsection{Classification}
On maximise la \textbf{pureté} des nœuds fils. Critères courants :
\begin{itemize}
    \item \textbf{Erreur de classification :} $1 - \max_k \hat{p}_{mk}$
    \item \textbf{Indice de Gini :} $\displaystyle G_m = \sum_{k \neq k'} \hat{p}_{mk} \hat{p}_{mk'} = 1 - \sum_k \hat{p}_{mk}^2$
    \item \textbf{Entropie croisée :} $\displaystyle D_m = -\sum_k \hat{p}_{mk} \log \hat{p}_{mk}$
\end{itemize}

Le \textbf{Gini} est le plus utilisé (CART). La division optimale minimise la moyenne pondérée du critère :
\[
\min_{j,s} \left[ n_g \cdot G_g + n_d \cdot G_d \right]
\]

% ============================================================
\section{Élagage (pruning)}
% ============================================================

Un arbre complet (\textit{full-grown tree}) surajuste presque toujours. On utilise l'\textbf{élagage par coût-complexité} (\textit{cost-complexity pruning}) :

\[
\text{Critère}_{\alpha}(A) = \underbrace{\sum_{m=1}^{|A|} \sum_{i \in R_m} (y_i - \hat{y}_{R_m})^2}_{\text{Erreur d'apprentissage}} + \alpha \cdot |A|
\]

où $|A|$ est le nombre de feuilles de l'arbre $A$ et $\alpha \ge 0$ un paramètre de régularisation.

\begin{itemize}
    \item $\alpha = 0$ : arbre complet (max de feuilles).
    \item $\alpha$ grand : arbre simple (peu de feuilles).
    \item On choisit $\alpha$ par \textbf{validation croisée} (souvent 10-fold CV).
\end{itemize}

% ============================================================
\section{Exemple numérique pas à pas}
% ============================================================

\subsection{Données}
Considérons $n = 6$ individus avec une variable explicative $X$ et une réponse $Y$ (régression) :

\begin{center}
\begin{tabular}{c|c}
$X$ & $Y$ \\
\hline
1 & 2 \\
2 & 3 \\
3 & 8 \\
4 & 9 \\
5 & 4 \\
6 & 5 \\
\end{tabular}
\end{center}

\subsection{Recherche de la première division}

On teste chaque seuil $s$ possible (entre 1 et 6).

\paragraph{Seuil $s = 1$ (X < 1 vs X $\ge$ 1) :}
\begin{itemize}
    \item Nœud gauche : $\emptyset$ (impossible — nœud vide)
    \item Nœud droit : tous les individus
\end{itemize}
Ce seuil n'est pas admissible (nœud vide).

\paragraph{Seuil $s = 2$ (X < 2 vs X $\ge$ 2) :}
\begin{itemize}
    \item Gauche : $\{1\}$, $\bar{y}_g = 2$, RSS$_g = 0$
    \item Droite : $\{2,3,4,5,6\}$, $\bar{y}_d = \frac{3+8+9+4+5}{5} = 5{,}8$, RSS$_d = (3-5{,}8)^2 + (8-5{,}8)^2 + (9-5{,}8)^2 + (4-5{,}8)^2 + (5-5{,}8)^2$
    \item RSS$_d = 7{,}84 + 4{,}84 + 10{,}24 + 3{,}24 + 0{,}64 = 26{,}8$
    \item RSS total = $0 + 26{,}8 = 26{,}8$
\end{itemize}

\paragraph{Seuil $s = 3$ (X < 3 vs X $\ge$ 3) :}
\begin{itemize}
    \item Gauche : $\{1,2\}$, $\bar{y}_g = 2{,}5$, RSS$_g = (2-2{,}5)^2 + (3-2{,}5)^2 = 0{,}5$
    \item Droite : $\{3,4,5,6\}$, $\bar{y}_d = \frac{8+9+4+5}{4} = 6{,}5$, RSS$_d = (8-6{,}5)^2 + (9-6{,}5)^2 + (4-6{,}5)^2 + (5-6{,}5)^2$
    \item RSS$_d = 2{,}25 + 6{,}25 + 6{,}25 + 2{,}25 = 17$
    \item RSS total = $0{,}5 + 17 = 17{,}5$
\end{itemize}

\paragraph{Seuil $s = 4$ (X < 4 vs X $\ge$ 4) :}
\begin{itemize}
    \item Gauche : $\{1,2,3\}$, $\bar{y}_g = \frac{2+3+8}{3} \approx 4{,}33$, RSS$_g \approx 20{,}67$
    \item Droite : $\{4,5,6\}$, $\bar{y}_d = \frac{9+4+5}{3} = 6$, RSS$_d = 14$
    \item RSS total $\approx 34{,}67$
\end{itemize}

\paragraph{Seuil $s = 5$ (X < 5 vs X $\ge$ 5) :}
\begin{itemize}
    \item Gauche : $\{1,2,3,4\}$, $\bar{y}_g = 5{,}5$, RSS$_g = 29$
    \item Droite : $\{5,6\}$, $\bar{y}_d = 4{,}5$, RSS$_d = 0{,}5$
    \item RSS total = $29{,}5$
\end{itemize}

\paragraph{Bilan :} Le seuil optimal est $s = 3$ (RSS minimal = 17{,}5). La première division est donc $X < 3$ vs $X \ge 3$.

\subsection{Division suivante}

On répète le processus sur chaque nœud fils.

\paragraph{Nœud gauche (individus 1, 2) :} $\bar{y} = 2{,}5$, variance intra = 0{,}25. Pas de gain significatif à diviser (trop petit). On s'arrête.

\paragraph{Nœud droit (individus 3, 4, 5, 6) :} On teste $s = 5$ (seuil optimal après calculs).
\begin{itemize}
    \item Gauche : $\{3,4\}$, $\bar{y}=8{,}5$, RSS = 0{,}5
    \item Droite : $\{5,6\}$, $\bar{y}=4{,}5$, RSS = 0{,}5
    \item RSS total = 1
\end{itemize}
Division retenue : $X < 5$ vs $X \ge 5$.

\subsection{Arbre final}

\begin{center}
\begin{tikzpicture}[
    node distance=1.8cm and 2.5cm,
    dec/.style={rectangle, draw, rounded corners, minimum width=2cm, minimum height=0.7cm, align=center, font=\bfseries, fill=blue!10},
    leaf/.style={rectangle, draw, minimum width=2.2cm, minimum height=0.7cm, align=center, font=\normalsize, fill=green!10},
    edge/.style={draw, -{Stealth[scale=1.2]}, thick},
    sn/.style={font=\small, above, sloped}
]

\node[dec] (root) {$X < 3$};

\node[leaf, below left=of root] (left1) {$Y = 2{,}5$ \\ \small (ind. 1, 2)};

\node[dec, below right=of root] (right1) {$X < 5$};

\node[leaf, below left=of right1] (left2) {$Y = 8{,}5$ \\ \small (ind. 3, 4)};

\node[leaf, below right=of right1] (right2) {$Y = 4{,}5$ \\ \small (ind. 5, 6)};

\draw[edge] (root) -- node[sn]{oui} (left1);
\draw[edge] (root) -- node[sn]{non} (right1);
\draw[edge] (right1) -- node[sn]{oui} (left2);
\draw[edge] (right1) -- node[sn]{non} (right2);

\end{tikzpicture}
\end{center}

Prédictions :
\[
\hat{y} = 
\begin{cases}
2{,}5 & \text{si } X < 3 \\
8{,}5 & \text{si } 3 \le X < 5 \\
4{,}5 & \text{si } X \ge 5
\end{cases}
\]

% ============================================================
\section{Code Python}
% ============================================================

\begin{lstlisting}[language=Python, caption=Arbre de régression/classification avec scikit-learn]
from sklearn.tree import DecisionTreeRegressor, DecisionTreeClassifier
from sklearn.tree import plot_tree
from sklearn.model_selection import cross_val_score
import matplotlib.pyplot as plt
import pandas as pd
import numpy as np

# --- Arbre de regression ---
np.random.seed(42)
X = np.linspace(0, 10, 100).reshape(-1, 1)
y = np.sin(X).ravel() + np.random.randn(100) * 0.1

tree_reg = DecisionTreeRegressor(max_depth=3, min_samples_leaf=5)
tree_reg.fit(X, y)

# Visualisation
plt.figure(figsize=(12, 6))
plot_tree(tree_reg, filled=True, feature_names=['X'])
plt.title("Arbre de regression")
plt.show()

# Prediction
X_new = np.array([[5.5]])
print(f"Prediction : {tree_reg.predict(X_new)[0]:.3f}")
\end{lstlisting}

\begin{lstlisting}[language=Python, caption=Arbre de classification — Iris]
from sklearn.datasets import load_iris
from sklearn.tree import DecisionTreeClassifier
from sklearn.model_selection import train_test_split
from sklearn.metrics import confusion_matrix, classification_report

iris = load_iris()
X, y = iris.data, iris.target

# Separation train/test
X_train, X_test, y_train, y_test = train_test_split(
    X, y, test_size=0.3, random_state=42)

# Arbre (Gini par defaut, profondeur max 4)
tree_clf = DecisionTreeClassifier(
    max_depth=4, criterion='gini', min_samples_leaf=5)
tree_clf.fit(X_train, y_train)

# Predictions
y_pred = tree_clf.predict(X_test)
print(confusion_matrix(y_test, y_pred))
print(classification_report(y_test, y_pred,
      target_names=iris.target_names))

# Importance des variables
for name, imp in zip(iris.feature_names, tree_clf.feature_importances_):
    print(f"{name}: {imp:.3f}")
\end{lstlisting}

\begin{lstlisting}[language=Python, caption=Élagage par coût-complexité]
# Chemin d'elagage (cost complexity pruning path)
path = tree_clf.cost_complexity_pruning_path(X_train, y_train)
ccp_alphas, impurities = path.ccp_alphas, path.impurities

# Selection par validation croisee
from sklearn.model_selection import GridSearchCV

param_grid = {'ccp_alpha': ccp_alphas}
grid = GridSearchCV(
    DecisionTreeClassifier(random_state=42),
    param_grid, cv=5
)
grid.fit(X_train, y_train)

print(f"Meilleur alpha : {grid.best_params_['ccp_alpha']:.4f}")
print(f"Meilleur score : {grid.best_score_:.3f}")

# Arbre elague optimal
best_tree = grid.best_estimator_
print(f"Nombre de feuilles : {best_tree.get_n_leaves()}")
\end{lstlisting}

\newpage
% ============================================================
\section{Code R}
% ============================================================

\begin{lstlisting}[language=R, caption=Arbre de classification avec rpart]
# install.packages("rpart")
# install.packages("rpart.plot")
library(rpart)
library(rpart.plot)

# Arbre de classification sur iris
data(iris)
set.seed(42)

# Separation train/test
idx <- sample(1:nrow(iris), 0.7 * nrow(iris))
train <- iris[idx, ]
test <- iris[-idx, ]

# Modele
arbre <- rpart(Species ~ ., data = train,
               method = "class",
               control = rpart.control(minsplit = 5,
                                       maxdepth = 4,
                                       cp = 0.01))

# Visualisation
rpart.plot(arbre, type = 2, extra = 104,
           main = "Arbre de classification (Iris)")

# Predictions
pred <- predict(arbre, test, type = "class")
table(pred, test$Species)

# Elagage par cp (complexity parameter)
print(arbre$cptable)  # table cout-complexite
best_cp <- arbre$cptable[which.min(arbre$cptable[, "xerror"]), "CP"]
arbre_elague <- prune(arbre, cp = best_cp)
rpart.plot(arbre_elague, type = 2, extra = 104,
           main = "Arbre elague optimal")
\end{lstlisting}

\begin{lstlisting}[language=R, caption=Arbre de régression]
# Arbre de regression
X <- seq(0, 10, length.out = 100)
y <- sin(X) + rnorm(100, sd = 0.1)
data_reg <- data.frame(X = X, Y = y)

tree_reg <- rpart(Y ~ X, data = data_reg,
                  method = "anova",
                  control = rpart.control(minsplit = 5,
                                          cp = 0.001))
rpart.plot(tree_reg, type = 2, main = "Arbre de regression")

# Prediction
predict(tree_reg, newdata = data.frame(X = 5.5))
\end{lstlisting}

\newpage
% ============================================================
\section*{QCM — Vérifiez vos connaissances}
% ============================================================

\medskip
\textit{Pour chaque question, choisissez la bonne réponse. Les réponses sont en page \pageref{reponses-tree}.}

\begin{enumerate}

    \item CART est l'acronyme de :
    \begin{enumerate}
        \item Classification And Regression Trees
        \item Categorical Automatic Rule Testing
        \item Clustering And Recursive Trees
        \item Continuous Asymptotic Regression Trees
    \end{enumerate}

    \item La construction d'un arbre CART procède par :
    \begin{enumerate}
        \item divisions aléatoires de l'espace
        \item partitionnement récursif binaire descendant
        \item fusion ascendante de régions
        \item optimisation globale par programmation dynamique
    \end{enumerate}

    \item En régression, le critère de division est :
    \begin{enumerate}
        \item l'indice de Gini
        \item l'entropie croisée
        \item la somme des carrés résiduels (RSS)
        \item le R²
    \end{enumerate}

    \item En classification, le critère de division le plus utilisé dans CART est :
    \begin{enumerate}
        \item l'erreur de classification
        \item l'indice de Gini
        \item l'entropie croisée
        \item le critère d'Akaike (AIC)
    \end{enumerate}

    \item L'élagage par coût-complexité vise à :
    \begin{enumerate}
        \item augmenter le nombre de feuilles
        \item réduire la variance en pénalisant la complexité
        \item accélérer l'apprentissage
        \item équilibrer les classes
    \end{enumerate}

    \item Dans le critère coût-complexité $\text{Err}(A) + \alpha|A|$, $|A|$ représente :
    \begin{enumerate}
        \item le nombre de nœuds internes
        \item le nombre de feuilles
        \item la profondeur de l'arbre
        \item le nombre de variables
    \end{enumerate}

    \item Pour choisir le paramètre d'élagage $\alpha$, on utilise :
    \begin{enumerate}
        \item un test de Student
        \item la validation croisée
        \item le critère AIC
        \item le test du chi-deux
    \end{enumerate}

    \item Les arbres CART sont dits \textbf{non-paramétriques} car :
    \begin{enumerate}
        \item ils n'ont pas de paramètres
        \item ils ne font pas d'hypothèse sur la distribution des données
        \item ils utilisent des paramètres non linéaires
        \item ils ne nécessitent pas d'apprentissage
    \end{enumerate}

    \item Un inconvénient majeur des arbres CART est :
    \begin{enumerate}
        \item leur forte variance (instabilité)
        \item leur incapacité à gérer les variables qualitatives
        \item leur temps d'apprentissage très long
        \item leur manque d'interprétabilité
    \end{enumerate}

    \item Quel package R implémente les arbres CART ?
    \begin{enumerate}
        \item randomForest
        \item rpart
        \item ggplot2
        \item caret
    \end{enumerate}

\end{enumerate}

\newpage
\label{reponses-tree}

\begin{center}
{\Large\textbf{Réponses au QCM}}
\end{center}

\medskip

\begin{tabular}{cp{2cm}p{9cm}}
\toprule
\textbf{N°} & \textbf{Réponse} & \textbf{Explication} \\
\midrule
1 & \textbf{a} & CART = Classification And Regression Trees (Breiman et al., 1984). \\
2 & \textbf{b} & L'algorithme divise récursivement chaque nœud en deux, de la racine vers les feuilles. \\
3 & \textbf{c} & En régression, on minimise la somme des carrés résiduels (RSS) dans les nœuds fils. \\
4 & \textbf{b} & L'indice de Gini $1 - \sum \hat{p}_{mk}^2$ est le critère par défaut de CART pour la classification. \\
5 & \textbf{b} & L'élagage réduit la variance en pénalisant un grand nombre de feuilles (contrôle du surajustement). \\
6 & \textbf{b} & $|A|$ est le nombre de feuilles (nœuds terminaux) de l'arbre. \\
7 & \textbf{b} & On choisit $\alpha$ par validation croisée (10-fold CV typiquement). \\
8 & \textbf{b} & Aucune hypothèse sur la distribution des données n'est nécessaire (contrairement à la régression linéaire). \\
9 & \textbf{a} & Un petit changement dans les données peut complètement modifier l'arbre (forte variance). \\
10 & \textbf{b} & Le package \texttt{rpart} implémente les arbres CART en R. \\
\bottomrule
\end{tabular}

\vspace{0.5cm}

\begin{center}
\rule{0.6\textwidth}{0.4pt}
\vspace{0.3cm}
\textit{Fin de la fiche.}
\end{center}

\end{document}
